\documentclass[../../../../main.tex]{subfiles}
\begin{document}

\begin{enumerate}
  \item$\triangle OP_0P_1$が与えられている．
辺$P_0O$上に点$P_2,P_4,\cdots,P_{2n},\cdots$を，
また辺$P_1O$上に点$P_3,P_5,\cdots,P_{2n+1},\cdots$を適当にとって，
次の三条件(イ)〜(ハ)がみたされるようにしうるための，
$\triangle OP_0P_1$についての条件を求めよ．
\begin{enumerate}
  \item$P_0,P_2,P_4,\cdots,P_{2n},\cdots$はこの順に線分$P_0O$上にならび，しかも$n$が大きくなるに従って，$O$に近づく．
  \item$P_1,P_3,P_5,\cdots,P_{2n+1},\cdots$はこの順に線分$P_1O$にならんでいる．
  \item$\triangle P_{n-1}P_nP_{n+1} \text{∽} \triangle P_nP_{n+1}P_{n+2}$が，
$n$のどの自然数値に対しても成立する．
(ここでいう相似は，
左辺の三角形の$P_{n-1}$，$P_n$，$P_{n+1}$が右辺の三角形の$P_n$，$P_{n+1}$，$P_{n+2}$に順次対応した相似を意味する．)
\end{enumerate}
  \item上のように点がとれたとき，$\triangle P_{n-1}P_nP_{n+1}$の面積を第$n$項$(n=1,2,\cdots)$とする数列は
\begin{enumerate}
  \itemいつでも等比数列である
  \item等比数列になる場合もあり，等比数列にならない場合もある
  \itemけっして等比数列にはならない
\end{enumerate}
のどれが正しいか，理由を付して答えよ．
\end{enumerate}

\end{document}
